% diagrams/firstwords/escenario.tex -- en las fiestas, Lucía, Martín y
% Dani leen su adivinanza en el escenario; Martín, con el papel y el
% micrófono.
\begin{tikzpicture}[dibujofw]
  % las cortinas y el escenario
  \filldraw[fill=white] (-7.6,0) -- (-7.6,10.0) -- (-4.6,10.0) .. controls (-5.6,7.0) and (-5.0,3.0) .. (-6.0,1.2) -- (-6.0,0) -- cycle;
  \filldraw[fill=white] (7.6,0) -- (7.6,10.0) -- (4.6,10.0) .. controls (5.6,7.0) and (5.0,3.0) .. (6.0,1.2) -- (6.0,0) -- cycle;
  \filldraw[fill=white] (-7.8,9.2) rectangle (7.8,10.2);
  \filldraw[fill=white] (-7.0,-1.4) rectangle (7.0,0.0);
  % Lucía, Martín (con su papel) y Dani
  \begin{scope}[shift={(-3.2,0)}] \ninaFW{Lucia}{Abajo}{Arriba} \end{scope}
  \begin{scope}[shift={(0,0)}, scale=0.9] \ninoFW{Martin}{Frente}{Frente}
    \filldraw[fill=white] (-1.1,3.0) -- (1.1,3.0) -- (1.0,5.0) -- (-1.0,5.0) -- cycle;
    \foreach \y in {3.5,4.0,4.5} {\draw (-0.7,\y) -- (0.7,\y);}
    \filldraw[fill=white] (-0.35,3.12) circle (0.27); \filldraw[fill=white] (0.35,3.12) circle (0.27);
  \end{scope}
  \begin{scope}[shift={(3.2,0)}, scale=0.85] \ninoFW{Dani}{Arriba}{Abajo} \end{scope}
  % el micrófono
  \draw (1.6,0) -- (1.6,4.4);
  \filldraw[fill=white] (1.3,0) -- (1.9,0) -- (1.6,0.3) -- cycle;
  \filldraw[fill=white] (1.6,4.7) ellipse (0.25 and 0.35);
\end{tikzpicture}
